% Part I of the guide "Lean through Artform Calculus": the basics of Lean,
% written 2026-10-09 at the author's request ("Add 10 pages in the beginning
% explaining the basics of lean in an accessible way"). Every listing is cut
% from ../LeanBasics.lean, which Lean 4.34.1 accepts; each answer of Lean
% quoted here is pinned there with #guard_msgs.
\chapter*{Part I. The basics of Lean}
\addcontentsline{toc}{chapter}{Part I. The basics of Lean}

Lean is a programming language in which programs and proofs are written in
one notation and checked by one program. A Lean file is a list of
declarations. Lean reads the file from top to bottom, and for each
declaration it either accepts it or says what is wrong with it. This part
teaches enough of the language to read Part~II, which writes the
constructions of Chapter~1 of the book in Lean. It assumes no experience of
programming.

Every example was checked by Lean~4.34.1. A comment after \lc{-{}-} at the
end of a line shows what Lean answers to that line. Where Lean refuses a line,
its message follows the listing, as Lean prints it.

\section*{Running Lean}
\addcontentsline{toc}{section}{Running Lean}

The quickest way to start needs nothing installed. The Lean playground at
\url{https://live.lean-lang.org} is an editor in the browser: type on the
left, and Lean answers on the right. For longer work, install the editor
VS~Code and its official Lean~4 extension. The extension installs Lean
itself and checks a file while it is being typed. Lean's answers appear in a
panel beside the text, called the \emph{infoview}, and an error is
underlined in red where it occurs. Lean also runs in a terminal:
\lc{lean Basics.lean} checks the file \lc{Basics.lean} and prints the answers
and the errors.

The examples of this guide are in the folder \lc{tutorial} of the
repository that holds it. Part~I's are in \lc{LeanBasics.lean} and
Part~II's in \lc{ArtformChapter1.lean}. Those files check every answer
quoted here, so an editor shows nothing while the answers are right. The
folder \lc{play} holds copies without the checks, which show every answer
in the infoview. To try them, copy \lc{play} to a folder outside the
repository and open it in VS~Code; inside the repository, the editor would
first set up the whole project, Mathlib included.

The examples use Lean's core language only. Most mathematicians who use Lean
also use Mathlib, a large library of mathematics written in Lean; nothing
here needs it.

\section*{Values and types}
\addcontentsline{toc}{section}{Values and types}

Every expression in Lean has a type. Two commands ask about an expression:
\li!#check! shows its type, and \li!#eval! computes its value.
\leanlisting{bvalues}
The colon is read \emph{has type}, as in the book: \li!2 + 3 : Nat! says that
$2+3$ is a natural number. \li!Nat! is the type of the natural numbers,
\li!String! of texts, \li!Nat × String! of pairs of a number and a text, and
\li!List Nat! of lists of numbers. A text is written between double quotes,
and \li!++! joins two texts.

Numbers need care. \li!Nat! has no negative numbers, so subtraction stops at
zero, and division of natural numbers drops the remainder, which \li!%!
gives. To compute in another type, write the expression with its type,
\li!(e : T)!. On the integers \li!Int! subtraction goes below zero, and on
the rationals \li!Rat! division is exact:
\leanlisting{bnumbers}
The last line is the sum $\frac14+\frac13=\frac7{12}$ of
\cref{fig:wheel-refinement}. Lean prints a rational with its type.

\section*{Definitions and functions}
\addcontentsline{toc}{section}{Definitions and functions}

\li!def! gives a name to a value. The name comes first; then the inputs, each
in parentheses with its type; then a colon and the type of the result; and
then \lc{:=} and the value.
\leanlisting{bdefs}
A function is applied by writing it before its input, without parentheses:
\li!double 21!. Parentheses only group, as in \li!double (double 3)!.
\li!#check! of a defined name shows its inputs and the type of its result.

A function is a value, and \li!fun! builds one without a name:
\li!fun n => 2 * n! is the function $n\mapsto2n$, of type \li!Nat → Nat!. A
function of two inputs takes them one at a time, so \li!add 2! is again a
function: the one that adds two to its input.
\leanlisting{bfunctions}
The first input of \li!twice! is a type, $\alpha$. The braces make it
\emph{implicit}: Lean works out $\alpha$ from the other inputs, here
\li!Nat!, and the caller does not write it. \li!twice double 3! doubles $3$
twice.

\section*{Definitions by cases, and recursion}
\addcontentsline{toc}{section}{Definitions by cases, and recursion}

A definition may be given case by case, with one line for each form the input
can take. A natural number is either $0$ or $n+1$ for a smaller number $n$,
and on the line for $n+1$ the definition may use itself at $n$:
\leanlisting{brecursion}
So \li!sumTo 4! is $4+3+2+1+0=10$, the ten points of the tetractys
(\cref{sec:interpretation}). Lean accepts such a definition only after two
checks: every case is covered, and every call of \li!sumTo! inside it is at a
smaller number, so that the computation ends. A definition that leaves a case
out is refused:
\leanlisting{bmissing}
\leanmessage{bmissing}
The message names the case that was left out: a number of the form $n+1$,
which Lean writes as \li!Nat.succ _!.

\section*{Inductive types}
\addcontentsline{toc}{section}{Inductive types}

A new type is made by listing the ways to build its values. Each line names a
\emph{constructor} and says what it takes. A colour is light or dark, and a
tree is either a leaf carrying a letter or a node joining two trees:
\leanlisting{binductive}
Every value of the type is built by its constructors, and only by them;
\cref{sec:rules-view} calls them the introduction rules. A function on a type
is defined by cases, one for each constructor, and on a node it may use
itself on the two smaller trees. The tree built from $a$, $b$ and $c$ below
has three leaves, as the construction trees of \cref{sec:formation-lines}
have one leaf for each atom.
\leanlisting{btreesize}
Two pieces of notation appear here. \li!Tree.size! is the function \li!size!
belonging to the type \li!Tree!, and for a tree \li!t!, \li!t.size! means
\li!Tree.size t!. Where Lean already knows that a tree is meant, \li!.leaf!
means \li!Tree.leaf!. The natural numbers are an inductive type too, with the
two constructors \li!zero! and \li!succ!, and the cases \li!0! and \li!n + 1!
of \li!sumTo! are these two.

\section*{Structures}
\addcontentsline{toc}{section}{Structures}

A \emph{structure} is a type of records with named fields. A value gives its
fields either by name or in order between angle brackets, and a field is read
with a dot:
\leanlisting{bstructure}
A structure is an inductive type with one constructor, which takes all the
fields, and \li!⟨2, 5⟩! applies that constructor.

\section*{Types that depend on values}
\addcontentsline{toc}{section}{Types that depend on values}

Types are values too: \li!Nat! has type \li!Type!. So a function may return a
type, and a type may depend on a number. \li!Fin n! is the type of the
numbers below $n$, the cells of the book's rod $R_n$. A value of it is a
number together with a proof that the number is below $n$, written between
angle brackets. The words \li!by decide! ask Lean to find the proof itself,
by computing.
\leanlisting{bfin}
The second line shows a trap. A bare numeral in \li!Fin 5! is read modulo
$5$, so \li!7! is accepted and means $2$, and nothing checks that $7$ is not
a cell of $R_5$. Part~II writes cells with their proofs for this reason.

A function whose output type depends on its input is a \emph{dependent
function}. \li!last n! is the largest number below $n+1$, and its type
\li!Fin (n + 1)! changes with $n$. Its proof, \li!by omega!, asks Lean to show
$n<n+1$ by arithmetic.
\leanlisting{bdependent}
The type of \li!last! is written \li!(n : Nat) → Fin (n + 1)!: it is the
book's $\AFDepFun_{n:\mathbb N}\mathsf{Fin}(n+1)$. In a \emph{dependent pair}
the type of the second entry depends on the first: a value of
\li!(n : Nat) × Fin n! is a number $n$ together with a number below it, the
book's $\AFDepPair_{n:\mathbb N}\mathsf{Fin}(n)$. The pair \li!p! is $3$ with
$1$.

\section*{Propositions are types}
\addcontentsline{toc}{section}{Propositions are types}

A proposition is a type as well, of the sort \li!Prop!, and a proof of it is a
value of that type. The equation \li!2 + 3 = 5! is a proposition, and
\li!rfl! is a proof of it, because both sides compute to $5$. \li!example!
states a proposition and gives its proof; \li!theorem! does the same and
gives the proof a name.
\leanlisting{bprop}
A false proposition has no proof. Asked to decide one by computing, Lean
answers that it is false:
\leanlisting{bfalse}
\leanmessage{bfalse}

The logical connectives are types built from types. A proof of an
implication $p\to q$ is a function that turns any proof of $p$ into a proof of
$q$. A proof of $p\wedge q$ is a pair of proofs. A proof of
$\exists x,\,P(x)$ is a pair: a witness $x$ together with a proof of $P(x)$. A
proof of $\forall x,\,P(x)$ is a dependent function, giving a proof of $P(x)$
for every $x$. \li!False! is the proposition with no proof, and $\neg p$ means
$p\to$\,\li!False!.
\leanlisting{bconnectives}
The first line proves $p\to p$ by the function that returns its input. The
second pairs two proofs. The third gives the witness $3$ with the proof that
$3\cdot3=9$. A proof of an existence statement is a witness with its
evidence, the shape of every witness in the book.

\section*{Tactics}
\addcontentsline{toc}{section}{Tactics}

A long proof is easier to build in steps. After \li!by! come \emph{tactics}:
commands that each change the \emph{goal}, the statement still to be proved,
until nothing is left to prove. In the editor the infoview shows the goal
after each step. These are the tactics this text uses:

\begin{center}
\small
\renewcommand{\arraystretch}{1.15}
\begin{tabularx}{\textwidth}{@{}>{\raggedright\arraybackslash}p{.17\textwidth} >{\raggedright\arraybackslash}X@{}}
\toprule
Tactic & What it does\\ \midrule
\lc{rfl} & proves an equation whose two sides compute to one value\\
\lc{decide} & checks a statement about finitely many cases by computing each\\
\lc{omega} & proves facts of integer arithmetic: $+$, $-$, $\le$, $<$, and products, quotients and remainders by fixed numbers\\
\lc{intro h} & begins a proof of $p\to q$ by assuming $p$ and naming its proof \lc{h}\\
\lc{exact e} & finishes the goal with the proof \lc{e}\\
\lc{cases}, \lc{induction} & split the goal by the constructors of a value\\
\lc{simp} & rewrites the goal with definitions and known equations\\
\lc{grind} & searches for a proof by combining the facts at hand\\
\bottomrule
\end{tabularx}
\end{center}

\leanlisting{btactics}
The first proof assumes $p$ and $p\to q$, then applies the second to the
first. The second is arithmetic. The third checks, for the ten numbers below
$10$, that a square leaves remainder $0$ or $1$ on division by $4$, never $2$.

Induction proves a statement for every natural number from two cases, $0$
and $n+1$; in the second case the statement for $n$ may be used. Here it
proves Gauss's formula $2(1+2+\cdots+n)=n(n+1)$:
\leanlisting{bgauss}
In the case \li!succ n! the name \li!ih! stands for the formula at $n$.
\li!simp only [sumTo]! unfolds \li!sumTo! once, at $n+1$, and \li!grind!
finishes the algebra.

\section*{Arguments that Lean finds}
\addcontentsline{toc}{section}{Arguments that Lean finds}

An input in square brackets is supplied by Lean, not by the caller.
\li![Add α]! asks for an addition on the type $\alpha$. Lean has one for
\li!Nat!, \li!Int! and \li!Rat!, so one definition serves all three:
\leanlisting{binstances}
Such an input is called an \emph{instance}, and the family of instances it
belongs to a \emph{class}. \li!decide! works in the same way: it asks for an
instance that says how to decide the statement, and Lean comes with instances
for numbers, for the types \li!Fin n!, and for statements built from them.
Part~II adds one class of its own, \li!Listed!, for the finite types of
Chapter~1.

\section*{Reading Lean's messages}
\addcontentsline{toc}{section}{Reading Lean's messages}

The commonest message is a type mismatch. It shows an expression, the type
it has, and the type Lean expected there:
\leanlisting{bmismatch}
\leanmessage{bmismatch}
Read it from the bottom. Lean expected a \li!Nat! on the right of \li!=!,
because the left side, \li!double 2!, is a \li!Nat!, and it found a
\li!String!. A misspelt name is reported as unknown:
\leanlisting{bunknown}
\leanmessage{bunknown}
When a message is hard to read, \li!#check! on the parts of the expression
usually shows which type is not the one meant.

\begin{literatureconnection}[title={Further reading}]
Avigad, de~Moura, Kong and Ullrich teach Lean from its types and tactics, in
the order of this part and much further \cite{AvigadTPIL}. Christiansen
teaches Lean as a programming language \cite{ChristiansenFPIL}. Avigad and
Massot teach mathematics in Lean with Mathlib \cite{AvigadMassotMIL}.
\end{literatureconnection}
